\begin{claim}{C.29}{Proved}{Circle centralizers and the winding of a plane loop}
For every subgroup $H\leq\operatorname{Homeo}(S^1)$, put
$C=C_{\operatorname{Homeo}^+(S^1)}(H)$, with the compact-open topology.
The inclusion induces an injection
\[
 \pi_1(C,\id)\longrightarrow
 \pi_1(\operatorname{Homeo}^+(S^1),\id)=\Z,              \tag{C.29.1}
\]
where the integer is the winding of the image of one point. Its image
lies in $\Z^H$, with $H$ acting on $\Z$ by its orientation character.
A based loop of winding zero has an explicit contraction in $C$ that
is natural under all homeomorphisms between circles.
In particular, a based null-homotopic loop in
$\operatorname{Homeo}^+(\R^2)$ preserving a Jordan curve restricts to
a winding-zero loop on that curve. Its restriction therefore has this
natural contraction against every commuting circle homeomorphism.
\end{claim}
\begin{proof}
Write the circle as $\R/\Z$, and set
\[
 \widetilde A=\{F\in\operatorname{Homeo}^+(\R):
                        F(x+1)=F(x)+1\}.
\]
The quotient map $\widetilde A\to\operatorname{Homeo}^+(S^1)$ is a
covering with deck transformations $F\mapsto F+n$, $n\in\Z$.
Indeed, requiring $F(0)$ to lie in a fixed interval of length less
than one gives continuous unique lifts over the corresponding
evaluation neighborhood. The homotopy
$(r,F)\mapsto(1-r)F+r\id$ contracts $\widetilde A$: strict
increase, the translation equation, and properness hold at every
parameter. Thus the target fundamental group in (C.29.1) is precisely
$\Z$. For a based loop $f_t$, its unique lift $F_t$ starting at
$\id$ ends at $x\mapsto x+n$, where $n$ is its point winding.

Suppose that $f_t\in C$ and $n=0$. Every lift $\widetilde h$ of an
element $h\in H$, including an orientation-reversing one, commutes
with every $F_t$. The commutator projects to the identity on the
circle, so it is an integer translation depending continuously on
$t$; it is the identity at $t=0$, and hence at all $t$.
Consequently $F_t$ is a based loop in the line centralizer of the
entire lift group of $H$. This group also contains the integer
translations. Apply C.27, explicitly
\[
 \begin{split}
 m_r(x)&=\min_{0\leq u\leq r}F_u(x),\qquad
 M_r(x)=\max_{0\leq u\leq r}F_u(x),\\
 Q_{s,t}(x)&=\operatorname{med}
       \bigl(m_{1-s}(x),F_t(x),M_{1-s}(x)\bigr).
 \end{split}                                                   \tag{C.29.2}
\]
Every $Q_{s,t}$ commutes with all lifts and with $x\mapsto x+1$,
so (C.29.2) descends to a contraction through elements of $C$.
This proves injectivity in (C.29.1).

The winding of a conjugated loop $h f_t h^{-1}$ is the winding of
$f_t$ multiplied by the orientation sign of $h$. Since these loops
agree for $h\in H$, the image is contained in $\Z^H$. For naturality,
let $q$ be a homeomorphism between two circles and choose any lift
$\widetilde q$ to their universal lines. The based lift of the
conjugated loop is $\widetilde q F_t\widetilde q^{-1}$.
An increasing $\widetilde q$ carries the two extrema in (C.29.2)
to the corresponding extrema; a decreasing one interchanges them.
Every line homeomorphism preserves the median of three points.
Thus (C.29.2) is carried to exactly the same formula for the
conjugated loop. Changing the lift of $q$ by an integer translation
does not change the descended contraction.

It remains to prove the stated plane winding comparison. Conjugate
by a Schoenflies homeomorphism so that the preserved Jordan curve is
the unit circle. Choose two distinct points $a,b$ on it. Evaluation
and the difference direction give the continuous maps
\[
 \operatorname{Homeo}^+(\R^2)
 \longrightarrow\operatorname{Conf}_2(\R^2)
 \longrightarrow S^1,
 \qquad
 f\longmapsto(f(a),f(b))
 \longmapsto\frac{f(b)-f(a)}{|f(b)-f(a)|}.                 \tag{C.29.3}
\]
For the given circle-preserving loop choose continuous angular lifts
$\alpha(t),\beta(t)$ with
$\alpha(t)<\beta(t)<\alpha(t)+2\pi$.
They both acquire the increment $2\pi n$: their difference remains
in $(0,2\pi)$ and returns to its initial value modulo $2\pi$.
The identity
\[
 e^{i\beta(t)}-e^{i\alpha(t)}
 =2i\,e^{i(\alpha(t)+\beta(t))/2}
            \sin\frac{\beta(t)-\alpha(t)}2
\]
shows that the winding of (C.29.3) is exactly $n$.
A null-homotopy of the ambient plane loop maps under (C.29.3) to a
null-homotopy in $S^1$, forcing $n=0$. This proves the exact
comparison and the last assertion.
\end{proof}

\paragraph{Application to C.12.}
The canonical lift of a hyperbolic centralizer loop is null-homotopic
as a loop of plane homeomorphisms: it is the image under the continuous
canonical lift homomorphism of a loop in the contractible $A^0$.
Every Jordan-curve orbit in C.28 is preserved throughout this loop.
Its restriction therefore has winding zero by (C.29.3), and (C.29.2)
contracts it naturally against its stabilizer in the full operator
lift group. Commuting operators carrying one circle orbit to another
carry these contractions to each other. This gives a circle-orbit
contraction with the exact required equivariance. C.34 subsequently
proves joint continuity, including inverses, on the full nonopen
locus. C.35 shows that this particular prescription need not admit
an equivariant extension over an adjacent open orbit, and C.36
constructs compatible alternatives on finite unions of closures
with two-dimensional proper operator image. The remaining global
compatibility requirements are stated in C.12, which remains open.
